\documentclass{article}
\usepackage{graphicx} % Required for inserting images
\usepackage{xurl}
\usepackage[T1]{fontenc}
\usepackage[czech]{babel}

\title{Implementace simulací}
\author{Lukáš Plachý}
\date{\today}

\begin{document}

\maketitle

\section{Implementace simulací}

Tato část popisuje praktickou implementaci čtyř navazujících modelů v prostředí
NetLogo. Všechny modely sdílejí stejnou základní strukturu: inicializaci
populace v proceduře \texttt{setup}, jeden simulační krok v proceduře
\texttt{go}, lokální přenos infekce v proceduře \texttt{infect}, změny stavů
v proceduře \texttt{change-state} a náhodný pohyb agentů v proceduře
\texttt{move}. Díky tomu lze výsledky jednotlivých variant porovnávat bez toho,
aby se současně měnila celá architektura modelu.

Modely jsou uloženy v souborech \path{examples/virus1.nlogox} až
\path{examples/virus4.nlogox}. Každý další model zachovává podstatnou část
předchozí verze a přidává jeden typ mechanismu navíc. Tato návaznost je pro
práci důležitá: rozdíly ve výsledcích lze interpretovat jako dopad konkrétního
rozšíření, nikoli jako důsledek kompletně jiné implementace.

\subsection{Společný rámec modelů}

Populace je reprezentována agenty typu \texttt{turtle}. Každý agent se pohybuje
v diskrétním prostoru NetLogo a interaguje pouze s agenty na stejném políčku
(\texttt{turtles-here}). Přenos infekce je tedy lokální a závisí na prostorovém
promíchání populace.

Základní zdravotní stavy jsou kódovány barvou agenta:
\begin{itemize}
\item \textbf{bílá} -- vnímavý agent (\textit{susceptible}),
\item \textbf{žlutá} -- exponovaný agent v inkubační době,
\item \textbf{červená} -- infekční agent,
\item \textbf{zelená} -- zotavený agent,
\item \textbf{modrá} -- očkovaný agent v Simulaci 4.
\end{itemize}

Každý agent nese interní čítače a příznaky, které řídí přechody mezi stavy:
\texttt{hours} počítá dobu strávenou v aktuálním stavu, \texttt{speed} určuje
rychlost pohybu, \texttt{illness-duration} uchovává délku nemoci,
\texttt{incubation-hours} délku inkubace a \texttt{infectious?} rozlišuje, zda
agent skutečně šíří infekci. Ve třetí a čtvrté simulaci se přidává
\texttt{immunity-hours}; ve čtvrté simulaci navíc \texttt{vaccinated?} a
\texttt{high-risk?}.

Jeden tick odpovídá jedné hodině modelového času. Krátkodobé modely
\texttt{virus1} a \texttt{virus2} končí ve chvíli, kdy v populaci nezůstává žádný
žlutý ani červený agent, případně nejpozději po 8760 ticích. Dlouhodobé modely
\texttt{virus3} a \texttt{virus4} běží do pevného limitu 8760 ticků, protože
umožňují reinfekce a infekce v nich nemusí přirozeně vymizet.

\subsection{Simulace 1: základní SEIR}

První model implementuje základní agent-based variantu SEIR bez úmrtnosti,
kapacitního omezení a zdrojů. V proceduře \texttt{setup} se vytvoří
\texttt{num-people} agentů, všichni začínají jako vnímaví. Náhodně vybraná
skupina \texttt{init-infected} agentů je nastavena do exponovaného stavu.

Přenos probíhá v proceduře \texttt{infect}. Infekční agenti, tedy červení agenti
s příznakem \texttt{infectious?}, zkoušejí nakazit ostatní vnímavé agenty na
stejném políčku. Pravděpodobnost přenosu je dána parametrem
\texttt{transmission-prob}. Nově nakažený agent přechází do žlutého stavu,
dostává délku inkubace \path{incubation-days * 24} a délku nemoci náhodně
vzorkovanou z Poissonova rozdělení kolem \path{mean-illness-duration}.

Přechody stavů zajišťuje procedura \texttt{change-state}. Po uplynutí
inkubační doby se žlutý agent mění na červeného a začíná být infekční. Po
uplynutí doby nemoci se červený agent mění na zeleného a přestává infekci šířit.
Protože model neobsahuje úmrtnost ani ztrátu imunity, zotavení agenti v této
variantě zůstávají v zeleném stavu až do konce simulace.

\subsection{Simulace 2: úmrtnost a kapacita systému}

Druhá simulace rozšiřuje základní SEIR model o globální proměnnou
\texttt{deaths} a o kapacitní omezení reprezentované parametrem
\texttt{healthcare-capacity}. Struktura pohybu, infekce a inkubace zůstává
stejná jako v Simulaci 1.

Rozdíl nastává při ukončení nemoci. Když červený agent dosáhne konce své
\texttt{illness-duration}, model vypočítá efektivní pravděpodobnost úmrtí.
Základní hodnota je \texttt{fatality-rate}. Pokud je počet současně infekčních
agentů vyšší než \texttt{healthcare-capacity}, pravděpodobnost úmrtí se násobí
\texttt{overload-multiplier}. Výsledná hodnota je omezena horní hranicí
100~\%.

Pokud agent zemře, zvýší se \texttt{deaths} a agent je odstraněn příkazem
\texttt{die}. Pokud nezemře, přechází do zeleného stavu. Tato implementace
umožňuje sledovat, zda se při stejné epidemické dynamice projeví přetížení
systému zvýšenou úmrtností.

\subsection{Simulace 3: zdroje, dočasná imunita a reinfekce}

Třetí simulace přidává explicitní zdrojovou vrstvu a dlouhodobou dynamiku.
Globální proměnná \texttt{current-resources} reprezentuje aktuální dostupnou
zásobu zdrojů. Při inicializaci je nastavena na \path{resource-stock} a v
každém kroku se navyšuje o \texttt{resource-regeneration}, nejvýše však zpět na
hodnotu \path{resource-stock}.

Při dokončení nemoci se model pokusí agentovi přidělit léčbu. Pokud je
\texttt{current-resources} alespoň \texttt{treatment-cost}, zdroje se sníží
o cenu léčby a agent je označen jako léčený. Pokud zdroje nestačí, agent léčený
není a jeho efektivní úmrtnost se násobí parametrem
\texttt{resource-penalty}. Nadále zůstává aktivní i kapacitní mechanismus ze
Simulace 2, takže se mohou kombinovat dva typy zátěže: přetížení systému a
nedostatek zdrojů.

Třetí simulace zároveň zavádí dočasnou imunitu. Zotavení agenti zůstávají
zelení, ale v proceduře \texttt{move} se jim zvyšuje \texttt{immunity-hours}.
Při kontaktu s infekčním agentem mohou být znovu nakaženi. Pravděpodobnost
reinfekce roste s časem od zotavení až do horizontu 180 dní. Tím model přestává
popisovat jednu uzavřenou epidemickou vlnu a umožňuje vznik dlouhodobé
perzistence infekce.

\subsection{Simulace 4: vakcinace, rizikové skupiny a strategie}

Čtvrtá simulace přidává heterogenitu agentů a jednoduché intervenční strategie.
Každý agent může mít příznak \texttt{vaccinated?} a \texttt{high-risk?}. Podíl
očkovaných určuje parametr \texttt{vaccination-rate}; podíl rizikových agentů
určuje \path{high-risk-share}. Očkovaní agenti jsou v rozhraní značeni modře.

Vakcinace ovlivňuje dvě části modelu. Při přenosu infekce má očkovaný agent
sníženou pravděpodobnost nákazy na jednu osminu původní hodnoty. Pokud se
očkovaný agent přesto nakazí, má při rozhodování o úmrtí sníženou základní
pravděpodobnost úmrtí na jednu pětadvacetinu. Tím vakcinace působí jak proti
šíření, tak proti závažnosti průběhu.

Parametr \texttt{strategy-mode} určuje, kdo může získat léčbu, pokud jsou
dostupné zdroje:
\begin{itemize}
\item \texttt{0} -- léčba je přidělena každému agentovi, pokud je dostatek zdrojů,
\item \texttt{1} -- léčba je přidělena pouze rizikovým agentům,
\item \texttt{2} -- léčba je přidělena pouze neočkovaným agentům.
\end{itemize}

Tento mechanismus je záměrně jednoduchý. Neřeší fronty ani optimalizaci v čase,
ale umožňuje porovnat, jak různé prioritizační pravidlo mění celkový počet
úmrtí a finální počet infekčních agentů při stejné základní epidemické dynamice.

\subsection{Metriky a reportery}

Všechny modely exportují stejnou sadu reporterů s prefixem \texttt{m\_}. Díky
tomu lze pro baseline scénáře používat jednotný agregační skript a stejnou
strukturu výstupních tabulek.

\begin{itemize}
\item \texttt{m\_final\_susceptible} -- počet bílých agentů na konci běhu,
\item \texttt{m\_final\_exposed} -- počet žlutých agentů na konci běhu,
\item \texttt{m\_final\_infected} -- počet červených agentů na konci běhu,
\item \texttt{m\_final\_recovered} -- počet zelených agentů na konci běhu,
\item \texttt{m\_ever\_infected} -- odhad počtu agentů, kteří nebyli na konci vnímaví,
\item \texttt{m\_ticks} -- délka běhu v ticích,
\item \texttt{m\_deaths} -- počet úmrtí,
\item \texttt{m\_final\_resources} -- finální zásoba zdrojů, případně \texttt{-1}
u modelů bez zdrojů.
\end{itemize}

Je třeba rozlišovat mezi finální hodnotou a časovým maximem. Aktuální
BehaviorSpace exporty ukládají finální metriky, nikoli celý průběh časových
řad. Proto se v kapitole výsledků u dlouhodobých modelů interpretuje
\path{m_final_infected} jako indikátor přetrvávající infekce na konci běhu,
ne jako epidemický peak.

\subsection{Automatizace experimentů}

Experimenty jsou definovány přímo v souborech \texttt{.nlogox} pomocí
BehaviorSpace. Každý model obsahuje baseline experiment \texttt{baseline\_auto}
s deseti opakováními. Nad rámec baseline jsou definovány rozšířené experimenty:
\begin{itemize}
\item \texttt{ofat\_transmission\_auto} v Simulaci 1,
\item \texttt{capacity\_overload\_grid\_auto} v Simulaci 2,
\item \texttt{resources\_grid\_auto} v Simulaci 3,
\item \texttt{strategy\_grid\_auto} v Simulaci 4.
\end{itemize}

Spouštění je zabaleno do skriptů v adresáři \texttt{scripts}. Skript
\path{run_behaviorspace.*} spouští jeden konkrétní experiment v headless
režimu NetLogo. Skripty \path{run_all_baselines.*} a
\texttt{run\_extended\_experiments.*} spouštějí připravené sady experimentů a
po každém běhu volají \texttt{aggregate\_results.py}. Skript
\texttt{run\_all\_experiments.*} pouze spojuje obě sady dohromady, aby bylo
možné jedním příkazem vytvořit všechny raw CSV i agregované summary soubory.

Agregační skript počítá pro zvolené metriky počet hodnot, průměr, směrodatnou
odchylku a kvartily. U baseline scénářů agreguje všechny běhy dohromady. U
rozšířených experimentů používá \texttt{--group-by}, takže výsledky počítá
odděleně pro každou kombinaci měněných parametrů.

\subsection{Vazba implementace na cíle práce}

Implementace podporuje cíl replikační studie třemi způsoby. Zaprvé, jednotlivé
modely tvoří čitelnou řadu od jednoduchého SEIR modelu po variantu se zdroji,
vakcinací a strategiemi. Zadruhé, všechny modely exportují kompatibilní metriky,
což umožňuje přímé srovnání baseline scénářů. Zatřetí, automatizace přes
BehaviorSpace a skripty v adresáři \texttt{scripts} umožňuje opakované spuštění
experimentů bez ručního klikání v NetLogo GUI.

Zároveň je nutné zdůraznit omezení implementace. Kontakty jsou lokální a
zjednodušené, populace je syntetická a model není kalibrován na reálná
epidemiologická data. Výsledky proto slouží hlavně k porovnání mechanismů a
strategií v kontrolovaném agent-based prostředí, nikoli k predikci reálné
epidemie.



\end{document}
